\documentclass[10pt]{article}
\usepackage[T1]{fontenc}
\usepackage[utf8]{inputenc}
\usepackage[french]{babel}
\usepackage{amsmath,amssymb,mathrsfs,mathpazo}
\usepackage{geometry}
\geometry{a4paper,margin=2cm}
\usepackage{graphicx,array,tabularx,tabularray,colortbl,multirow,multicol,booktabs,enumitem,pifont,ulem,fancybox,diagbox,cancel}
\usepackage{pstricks,pstricks-add,pst-plot,pst-tree,pst-node,pst-text,pst-func,pst-math,pst-bspline,pst-3dplot,pst-eucl,pst-solides3d}
\let\pscalculate\undefined
\usepackage{pst-bezier}
\usepackage{tikz,pgfplots,tkz-tab}
\usetikzlibrary{arrows,arrows.meta,patterns,calc,positioning,trees,decorations.markings,shapes.geometric,shapes.misc,shapes.arrows,intersections,decorations.pathmorphing,decorations.pathreplacing,backgrounds,shadows,babel}
\pgfplotsset{compat=1.18}
\usepgfplotslibrary{fillbetween}
\usepackage[np]{numprint}
\usepackage{listings,makecell,pas-tableur,twoopt}
\let\barre\undefined
\usepackage[pstricks]{tablvar}
\definecolor{darkgreen}{RGB}{0,100,0}
\providecommand{\red}{\color{red}}
\providecommand{\blue}{\color{blue}}
\providecommand{\green}{\color{green}}
\providecommand{\black}{\color{black}}
\providecommand{\white}{\color{white}}
\providecommand{\gray}{\color{gray}}
\providecommand{\cyan}{\color{cyan}}
\providecommand{\magenta}{\color{magenta}}
\providecommand{\yellow}{\color{yellow}}
\providecommand{\violet}{\color{violet}}
\providecommand{\orange}{\color{orange}}
\providecommand{\euro}{\texteuro}
\providecommand{\EUR}{\texteuro}
\providecommand{\ieme}{\textsuperscript{e}}
\providecommand{\degres}{\ensuremath{^\circ}}
\providecommand{\degree}{\ensuremath{^\circ}}
\providecommand{\vv}[1]{\overrightarrow{#1}}
\providecommand{\Coord}[1]{#1}
\providecommand{\decosix}{}
\providecommand{\starredbullet}{\ensuremath{\bullet}}
\pagestyle{empty}
\setlength{\parindent}{0pt}
\usepackage{adjustbox}\DeclareUnicodeCharacter{2236}{:}\DeclareUnicodeCharacter{2190}{\ensuremath{\leftarrow}}
\providecommand{\R}{\mathbb{R}}
\providecommand{\N}{\mathbb{N}}
\providecommand{\Z}{\mathbb{Z}}
\providecommand{\Q}{\mathbb{Q}}
\providecommand{\D}{\mathbb{D}}
\providecommand{\C}{\mathbb{C}}
\providecommand{\e}{{\mathrm{e}}}
\providecommand{\vect}[1]{\overrightarrow{#1}}
\providecommand{\vectt}[1]{\overrightarrow{\mathrm{#1}}}
\providecommand{\vvt}[1]{\overrightarrow{\mathrm{#1}}}
\providecommand{\np}[2][]{{#2~#1}}
\providecommand{\vv}[1]{{\overrightarrow{#1}}}
\renewcommand{\d}{{\mathrm{d}}}
\providecommand{\degres}{{^\circ}}
\providecommand{\degre}{{^\circ}}
\providecommand{\blue}{\color{blue}}
\providecommand{\red}{\color{red}}
\providecommand{\enskip}{\quad}
\providecommand{\llbracket}{\lbrack\!\lbrack}
\providecommand{\rrbracket}{\rbrack\!\rbrack}
\providecommand{\euro}{{\text{€}}}
\providecommand{\no}{{\text{n°}}}
\providecommand{\interval}[2]{{[#1;#2]}}
\providecommand{\intervFO}[2]{{[#1;#2[}}
\providecommand{\dotfill}{\ldots}

\providecommand{\ds}{\displaystyle}
\providecommand{\textding}[1]{#1}
\providecommand{\Oij}{\(\mathrm{O};\vec{\imath},\vec{\jmath}\)}
\providecommand{\Oijk}{\(\mathrm{O};\vec{\imath},\vec{\jmath},\vec{k}\)}
\providecommand{\Ouv}{\(\mathrm{O};\vec{u},\vec{v}\)}
\providecommand{\intervFF}[2]{\ensuremath{[#1;#2]}}
\renewcommand{\euro}{\ifmmode\text{\texteuro}\else\texteuro\fi}
\DeclareUnicodeCharacter{00B0}{\ifmmode^{\circ}\else\textdegree\fi}
\DeclareUnicodeCharacter{20AC}{\ifmmode\text{\texteuro}\else\texteuro\fi}
% Énoncés : Éducation nationale. Transcription et corrigés : APMEP, Denis Vergès et contributeurs. Découpage et classement : CoopMaths, Sébastien Lozano, Éric Elter et contributeurs. Adaptation de la présentation : Maths Explorer.
% CC BY-SA : https://coopmaths.fr/www/about/
% https://coopmaths.fr/alea/static/bac/2024/tex/bac_2024_06_sujet1_suede_1_cor.tex
\begin{document}
\noindent\textit{APMEP / CoopMaths. CC BY-SA. Présentation adaptée par Maths Explorer.}\par\medskip
\textbf{Affirmation 1 : Fausse.}\quad
En effet :

Dérivons la fonction $f$ (qui est dérivable sur $\R$, en tant que somme d'une fonction polynôme, définie et dérivable sur $\R$, et d'une fonction composée d'une fonction linéaire suivie de la fonction exponentielle, toutes deux définies et dérivables sur $\R$).

$\forall x \in \R,\quad f'(x) = 2\e^{2x} - 6$.

Donc :\quad $\begin{aligned}[t] \forall x\in \R,\quad f'(x) - 2f(x)
	&	= \big(2\e^{2x} - 6\big) - 2\big(\e^{2x} - 6x + 1\big)\\
	&	= 2\e^{2x} - 6 - 2\e^{2x} + 12x -2 \quad\text{en distribuant le }- 2\\
	&	= 12x -8\\	\end{aligned}$

Ainsi, on a \quad $f' - 2f$\quad qui est la fonction \quad $x \longmapsto 12x - 8$,\quad qui est différente de la fonction \quad $x \longmapsto -6x + 1$,\quad donc $f$ n'est pas une solution de l'équation différentielle (E).

\medskip

\textbf{Affirmation 2 : Fausse}\quad 	En effet :

En posant la suite $\left(v_n\right)$, suite géométrique de premier terme $v_0 = 1$ et de raison $q = \dfrac{3}{4}$, on a donc :\quad
	$\forall n \in \N, \quad v_n = 1\times \left(\dfrac{3}{4}\right)^n = \left(\dfrac{3}{4}\right)^n$

Et donc  : \quad $\forall n \in \N,\quad u_n = v_0 + v_1 + v_2 + \dots +v_n$.

En utilisant la formule connue pour déterminer la somme des $(n+1)$ premiers termes d'une suite géométrique, on a :

$\forall n \in \N,\quad u_n = v_0 \times \dfrac{1 - q^{n+1}}{1 - q} = 1\times \dfrac{1 - \left(\dfrac{3}{4}\right)^{n+1}}{1 - \dfrac{3}{4}} = \dfrac{1 - \left(\dfrac{3}{4}\right)^{n+1}}{\dfrac{1}{4}} = 4\left(1 - \left(\dfrac{3}{4}\right)^{n+1}\right)$.

Comme on a :\quad $0< \dfrac{3}{4} < 1$, \quad on en déduit : \quad$\lim\limits_{n \to +\infty} \left(\dfrac{3}{4}\right)^{n+1} = 0$.

Puis, par limite de la somme, on a : \quad$\lim\limits_{n \to +\infty} 1 - \left(\dfrac{3}{4}\right)^{n+1} = 1$.

Et, enfin, par limite du produit :\quad$\lim\limits_{n \to +\infty} 4\left( 1-\left(\dfrac{3}{4}\right)^{n+1}\right) = 4$.

La limite de la suite est donc finie (et vaut 4).

\medskip

\textbf{Affirmation 3 : Fausse}.\quad En effet :

La fonction \texttt{suite()} part d'une variable \texttt{S} qui prend une valeur égale à 0, puis effectue une boucle \texttt{for}, qui va ajouter à ce 0 des puissances de $\dfrac{3}{4}$ (l'instruction \texttt{a**b} signifie $a^b$).

Il y a ici deux problèmes dans cette fonction pour que \texttt{suite(50)} renvoie la valeur $u_{50}$.

\begin{itemize}[label=$\bullet$]
	\item La première erreur, est qu'au lieu d'ajouter des puissances successives de $\dfrac{3}{4}$, on ajoute toujours la même puissance de $\dfrac{3}{4}$. L'appel \texttt{suite(50)} ajoute $\left(\dfrac{3}{4}\right)^{50}$ à chaque exécution de la boucle \texttt{for}, au lieu d'ajouter $\left(\dfrac{3}{4}\right)^{0}$, puis $\left(\dfrac{3}{4}\right)^{1}$, puis $\left(\dfrac{3}{4}\right)^{2}$ et ainsi de suite jusqu'à $\left(\dfrac{3}{4}\right)^{50}$.
	\item 	La seconde erreur est que le nombre $u_{50}$ est la somme des 51 premières puissances de $\dfrac{3}{4}$ (de la puissance 0 à la puissance 50), et la boucle \texttt{for} ici s'exécutera 50 fois seulement, et pas 51 fois.
	\end{itemize}

	Pour que l'appel \texttt{suite(50)} renvoie $u_{50}$, il faudrait une fonction comme celle-ci :

\begin{center}
\begin{ttfamily}
\begin{tabularx}{8cm}{c|X|}\cline{2-2}
	1&def suite(k):\\
	2&\quad S=0\\
	3&\quad for i in range({\red{k+1}}):\\
	4&\qquad S=S+(3/4)**{\red i}\\
	5&\quad return S\\ \cline{2-2}
\end{tabularx}
\end{ttfamily}
\end{center}

La ligne 3 modifiée permet d'exécuter la boucle \texttt{for} 51 fois (la variable \texttt{i} commencera à la valeur 0, puis augmentera d'unité en unité jusqu'à atteindre la valeur strictement inférieure à \texttt{k+1}, donc d'atteindre la valeur \texttt{k}).

La ligne 4 modifiée fait qu'à chaque exécution de la boucle for, on va ajouter une puissance différente de la raison, avec un exposant de plus en plus élevé.

\emph{Remarques :} l'instruction \texttt{range(k+1)} a exactement le même effet que \texttt{range(0,k+1)}.

On peut aussi trouver d'autres fonctions, avec un code différent qui auront le même effet. Il n'est pas nécessaire pour répondre à la question de donner une fonction correcte : pointer l'une des deux erreurs commises suffit à montrer que l'appel \texttt{suite(50)} ne produit pas l'effet escompté.

\medskip

\textbf{Affirmation 4 : Vraie.} \quad En effet :

Quel que soit le réel $a$, la fonction $f$ est définie et dérivable sur $]0~;~+\infty[ $, en tant que somme d'une fonction linéaire et du produit de la fonction ln par un réel, toutes fonctions définies et dérivables sur l'intervalle considéré.

Soit $a$ un réel :

On a :\quad $\forall x \in ]0 ; +\infty[, \quad f'(x) = a\times \dfrac{1}{x} - 2 = \dfrac{a}{x} - 2$.

Notamment, pour $x = 1 $\quad $f'(1) = \dfrac{a}{1} - 2 = a-2$.

La tangente à $C$ au point d'abscisse 1 est parallèle à l'axe des abscisses si et seulement si $f'(1) = 0$, c'est-à-dire si et seulement si $a = 2$.

Il existe donc une (unique) valeur de $a$ telle que la tangente à $C$ au point d'abscisse 1 est parallèle à l'axe des abscisses.

\end{document}
